\providecommand{\devnum}[1]{#1}

\section{Problem Model}
\label{sec:model}

\subsection{Jobs, Servers and Dispatch}

A finite set of $n$ jobs is served by $k$ identical servers. We say \emph{server} for what a
deployment calls an executor, a judging machine or a worker, and \emph{job} for a
submission, a function call or a test run. Job $i$ arrives at time $a_i \ge 0$ and needs
$C_i$ units of service, which we also call its \emph{cost} or its \emph{work}. Service is
non-preemptive. A server that starts a job keeps it until the job finishes or has run for
the time limit $L$, at which point the system stops it. Since $C_i$ is the work actually
executed, $0 < C_i \le L$, and a job stopped at the limit contributes exactly $L$. For a
stopped job, the waiting-time guarantee says nothing about the lost result. A \emph{task}
is the programming exercise a job attempts, so many jobs share one task.

Servers are work-conserving, which means that no server is idle while a job waits. Jobs are
numbered in arrival order, with ties broken by input order, and we write $j < i$ when $j$
comes first in this order, so $j < i$ implies $a_j \le a_i$. At a common instant,
completions are processed first, then arrivals, then dispatches. A \emph{policy} is any
rule that, whenever a server is free and a job waits, picks one waiting job and starts it.
Such a rule may be randomised, learned, priority-driven or driven by predictions of any
quality. We make no stochastic assumption about the arrival times, the costs or their
dependence, because the load we care about follows a calendar and is not stationary.

\subsection{What the Scheduler Knows}

A job's cost is hidden when it arrives, and the scheduler learns $C_i$ only when $i$
completes. At arrival it sees a feature vector $x_i$, built from the history of earlier
jobs, and a score $\hat C_i$ computed from $x_i$. This history must be causal. Write
$\mathrm{done}_j$ for the instant at which the result of job $j$ becomes readable, and
$\delta \ge 0$ for the delay before it reaches the store the features are read from. Then
\begin{equation}
  \text{information about } j \text{ may enter } x_i
  \quad\Longleftrightarrow\quad
  \mathrm{done}_j + \delta \le a_i .
  \label{eq:visibility}
\end{equation}
In a simulation driven by a log, $\mathrm{done}_j$ can be read in two ways. One of the two
ways takes the time at which the result appeared on the original platform, as the log
records it. However, in the simulated queue the job completes at a time that depends on the
policy being simulated. A feature that is
causal on the original platform's clock may therefore use a result that the simulated
queue has not produced yet. On the original clock we run the prediction study of
Section~\ref{sec:prediction}, and on the simulated one the headline scheduling results
(Section~\ref{sec:exp_visibility}). The guarantees of Section~\ref{sec:theory} hold under
either reading, since they assume nothing about the scores.

\subsection{Waiting Time, the FCFS Reference and the Service Guarantee}

Let $s^P_i$ be the start time of job $i$ under policy $P$ and $W_P[i] = s^P_i - a_i$ its
waiting time. FCFS always starts the earliest-arrived waiting job, and when several
servers become free at one instant it starts jobs in arrival order. We run FCFS on the same
arrivals, costs and $k$ as $P$, so that the two are compared job by job on one sample path.
The \emph{excess} of job $i$ is
\begin{equation}
  \mathrm{excess}_P[i] \;=\; W_P[i] - W_{\mathrm{FCFS}}[i] .
  \label{eq:excess}
\end{equation}
Writing $j \prec i$ when $j$ is started before $i$, the work that overtook $i$ and the
work that $i$ overtook are
\begin{equation}
  \mathrm{In}_i \;=\; \sum_{j > i,\; j \prec i} C_j ,
  \qquad
  \mathrm{Out}_i \;=\; \sum_{j < i,\; i \prec j} C_j .
  \label{eq:inout}
\end{equation}

A policy $P$ offers a \emph{service guarantee} $G$ if $\mathrm{excess}_P[i] \le G$ for every
job $i$, on every arrival sequence, for all costs bounded by $L$ and for every score,
adversarial scores included. An operator states $G$ in advance, and the guarantee is the
proved statement that every job stays within it. This guarantee
holds for each job on the realised sample path, not on average, and it assumes nothing about
the quality of the score. Packet schedulers state additive per-flow bounds of this kind
against a reference discipline \cite{parekh1993gps}. Batch schedulers measure the same
difference after the fact, against a fair start time \cite{sabin2004fairness}. However,
the scheduling-with-predictions literature compares with an optimal policy
\cite{purohit2018predictions,scully2022uniform}. We take FCFS as the reference because
these services run it by default, so the excess is a quantity an operator can be held
accountable for.

\subsection{Metrics}

The primary metric is the 99th percentile (p99) of the waiting time of the jobs that
arrive in the 24 hours before a deadline, written $\mathrm{p99}_{\mathrm{dl}}$. This metric
was fixed before any sealed part of a log was opened. With it we report the mean wait over
all jobs and the largest excess, which is what a guarantee bounds. We also report the
\emph{harm}, the largest excess among the jobs that FCFS would have started within 1 s, that
is, among jobs FCFS would not have delayed. A fourth quantity, the \emph{firing rate}, is
the share of dispatch decisions at which the guard of Section~\ref{sec:scheduling} overrides
the base policy, weighted by queue length. Finally, we report the worst wait among heavy
jobs, defined in Section~\ref{sec:prediction}.

Improvements are expressed as the fraction of the gap closed,
\begin{equation}
  \mathrm{gap\ closed}(P) \;=\;
  \frac{M_{\mathrm{FCFS}} - M_{P}}{M_{\mathrm{FCFS}} - M_{\mathrm{SJF}}} ,
  \label{eq:gapclosed}
\end{equation}
where $M_Q$ is $\mathrm{p99}_{\mathrm{dl}}$ under policy $Q$ and SJF orders the waiting
jobs by their true costs. SJF cannot be deployed, since costs are unknown at arrival. On
several non-preemptive servers SJF also minimises none of our metrics, so it is a reference
and not an optimum. A gap closed of 1 means that $P$ matches the p99 of SJF, and 0 that it
matches FCFS. Gap closed is not a percentage reduction in waiting time, and the tables
print the absolute waits beside it.

\subsection{Steady-State Models and This Load}

Stationary models give closed forms for the mean wait of each class of a non-preemptive
priority queue \cite[pp.~499--507]{harcholbalter2013}. The Schedule Ordered by Age-based
Priority (SOAP) framework analyses every policy that ranks jobs by a function of their size
class and age \cite{scully2018soap}, a family that contains ranking by predicted cost. Both
analyses need stationary arrivals, which deadline-driven load does not supply. On a synthetic Poisson
stream the priority formula agrees with simulation to \devnum{0.4\%}, but on the same jobs
at their recorded arrival times the measured mean wait is \devnum{10.7} to \devnum{13.7}
times what the formula predicts (Supplementary Section~S2.5). Thus, we use steady-state
results as motivation only, and state every guarantee on the realised sample path.
